\textbf{Phần 1. Trắc nghiệm lựa chọn}
\setcounterpageref{ex}{0}

\begin{ex} %[2H5V26]
Đường gấp khúc gồm $3$ đoạn thẳng có độ dài lần lượt là $15\text{ cm}, 24\text{ cm}, 18\text{ cm}$. Độ dài đường gấp khúc đó là
\choice
{$55\text{ cm}$}
{\True $57\text{ cm}$}
{$56\text{ cm}$}
{$58\text{ cm}$}
\loigiai{
$15 + 24 + 18 = 57\text{ (cm)}$.
\textbf{Đáp án của câu hỏi là B}
}
\end{ex}

\begin{ex} %[2H5V26]
Một hình tứ giác có độ dài các cạnh lần lượt là $12\text{ cm}, 15\text{ cm}, 18\text{ cm}, 20\text{ cm}$. Chu vi của hình tứ giác đó là
\choice
{$63\text{ cm}$}
{$64\text{ cm}$}
{$66\text{ cm}$}
{\True $65\text{ cm}$}
\loigiai{
$12 + 15 + 18 + 20 = 65\text{ (cm)}$.
\textbf{Đáp án của câu hỏi là D}
}
\end{ex}

\begin{ex} %[2H5V26]
Một đường gấp khúc gồm $4$ đoạn thẳng dài bằng nhau, mỗi đoạn dài $6\text{ cm}$. Chiều dài đường gấp khúc đó là
\choice
{\True $24\text{ cm}$}
{$20\text{ cm}$}
{$22\text{ cm}$}
{$26\text{ cm}$}
\loigiai{
$6 \times 4 = 24\text{ (cm)}$.
\textbf{Đáp án của câu hỏi là A}
}
\end{ex}

\begin{ex} %[2H5V26]
Đường gấp khúc $ABCD$ dài $62\text{ cm}$, trong đó đoạn $AB = 18\text{ cm}$, $BC = 25\text{ cm}$. Đoạn thẳng $CD$ dài bao nhiêu?
\choice
{$17\text{ cm}$}
{$18\text{ cm}$}
{\True $19\text{ cm}$}
{$20\text{ cm}$}
\loigiai{
$CD = 62 - (18 + 25) = 62 - 43 = 19\text{ (cm)}$.
\textbf{Đáp án của câu hỏi là C}
}
\end{ex}

\begin{ex} %[2H5V26]
Hình tứ giác có bao nhiêu đỉnh và bao nhiêu cạnh?
\choice
{$3$ đỉnh và $4$ cạnh}
{\True $4$ đỉnh và $4$ cạnh}
{$4$ đỉnh và $3$ cạnh}
{$5$ đỉnh và $5$ cạnh}
\loigiai{
Hình tứ giác luôn có đúng 4 đỉnh và 4 cạnh.
\textbf{Đáp án của câu hỏi là B}
}
\end{ex}

\begin{ex} %[2H5V26]
Uốn một sợi dây thép dài $40\text{ cm}$ thành một hình tứ giác có $4$ cạnh bằng nhau. Độ dài mỗi cạnh của hình tứ giác đó là
\choice
{$8\text{ cm}$}
{$9\text{ cm}$}
{$12\text{ cm}$}
{\True $10\text{ cm}$}
\loigiai{
$40 : 4 = 10\text{ (cm)}$.
\textbf{Đáp án của câu hỏi là D}
}
\end{ex}

\begin{ex} %[2H5V26]
Đường gấp khúc $MNPQ$ gồm $3$ đoạn thẳng. Biết $MN = 1\text{ dm } 2\text{ cm}$, $NP = 15\text{ cm}$, $PQ = 23\text{ cm}$. Chiều dài đường gấp khúc $MNPQ$ là
\choice
{\True $50\text{ cm}$}
{$48\text{ cm}$}
{$49\text{ cm}$}
{$52\text{ cm}$}
\loigiai{
Đổi $1\text{ dm } 2\text{ cm} = 12\text{ cm}$. Tổng chiều dài: $12 + 15 + 23 = 50\text{ (cm)}$.
\textbf{Đáp án của câu hỏi là A}
}
\end{ex}

\begin{ex} %[2H5V26]
Hình chữ nhật có phải là một hình tứ giác không?
\choice
{Không, hình chữ nhật không phải hình tứ giác}
{Chỉ đúng khi các cạnh bằng nhau}
{\True Có, hình chữ nhật là một hình tứ giác đặc biệt}
{Chưa thể khẳng định}
\loigiai{
Hình chữ nhật có 4 cạnh và 4 góc vuông nên là một hình tứ giác đặc biệt.
\textbf{Đáp án của câu hỏi là C}
}
\end{ex}

\begin{ex} %[2H5V26]
Một hình tam giác có chu vi $36\text{ cm}$, biết độ dài hai cạnh là $12\text{ cm}$ và $15\text{ cm}$. Độ dài cạnh thứ ba là
\choice
{$8\text{ cm}$}
{\True $9\text{ cm}$}
{$10\text{ cm}$}
{$11\text{ cm}$}
\loigiai{
Cạnh thứ ba: $36 - (12 + 15) = 36 - 27 = 9\text{ (cm)}$.
\textbf{Đáp án của câu hỏi là B}
}
\end{ex}

\begin{ex} %[2H5V26]
Đường gấp khúc $ABC$ có độ dài $28\text{ cm}$. Biết đoạn thẳng $AB$ dài hơn đoạn thẳng $BC$ là $4\text{ cm}$. Độ dài đoạn thẳng $AB$ là
\choice
{$14\text{ cm}$}
{$15\text{ cm}$}
{$18\text{ cm}$}
{\True $16\text{ cm}$}
\loigiai{
Đoạn $AB = (28 + 4) : 2 = 16\text{ (cm)}$.
\textbf{Đáp án của câu hỏi là D}
}
\end{ex}

\begin{ex} %[2H5V26]
Kẻ một đoạn thẳng chia hình tứ giác thành hai hình. Hai hình đó có thể là
\choice
{\True Hai hình tam giác}
{Hai hình tròn}
{Một hình tròn và một hình tam giác}
{Hai đoạn thẳng}
\loigiai{
Nối hai đỉnh đối diện của hình tứ giác ta được hai hình tam giác.
\textbf{Đáp án của câu hỏi là A}
}
\end{ex}

\begin{ex} %[2H5V26]
Cho đường gấp khúc khép kín gồm $4$ đoạn thẳng. Hình tạo bởi đường gấp khúc đó là
\choice
{Hình tam giác}
{Hình ngũ giác}
{\True Hình tứ giác}
{Đường cong kín}
\loigiai{
Đường gấp khúc khép kín 4 đoạn thẳng tạo thành một hình tứ giác.
\textbf{Đáp án của câu hỏi là C}
}
\end{ex}

\textbf{Phần 2. Trắc nghiệm đúng sai}
\setcounterpageref{ex}{0}

\begin{ex} %[2H5V26]
Xét tính đúng sai của các nhận định về đường gấp khúc và hình tứ giác:
\choiceTF
{\True Độ dài đường gấp khúc bằng tổng độ dài các đoạn thẳng tạo nên nó.}
{Đường gấp khúc gồm 3 đoạn thẳng thì có 4 điểm thẳng hàng.}
{\True Chu vi hình tứ giác bằng tổng độ dài 4 cạnh của hình tứ giác đó.}
{Hình tam giác có thể có 4 cạnh.}
\loigiai{
\begin{itemize}
    \item Ý a) Đúng: Đúng theo định nghĩa tính độ dài đường gấp khúc.
    \item Ý b) Sai: Các điểm trong đường gấp khúc không được cùng thẳng hàng.
    \item Ý c) Đúng: Đúng.
    \item Ý d) Sai: Hình tam giác chỉ có đúng 3 cạnh.
\end{itemize}
\begin{itemize}
    \item a) \textbf{Đúng}.
    \item b) \textbf{Sai}.
    \item c) \textbf{Đúng}.
    \item d) \textbf{Sai}.
\end{itemize}
}
\end{ex}

\begin{ex} %[2H5V26]
Một hình tứ giác MNPQ có MN = 15 cm, NP = 20 cm, PQ = 25 cm, QM = 18 cm:
\choiceTF
{\True Cạnh dài nhất của tứ giác là cạnh PQ.}
{\True Chu vi hình tứ giác là 78 cm.}
{Cạnh ngắn nhất của tứ giác là cạnh QM.}
{\True Tổng hai cạnh MN và PQ bằng 40 cm.}
\loigiai{
\begin{itemize}
    \item Ý a) Đúng: $25\text{ cm}$ là dài nhất. Đúng.
    \item Ý b) Đúng: $15 + 20 + 25 + 18 = 78\text{ cm}$.
    \item Ý c) Sai: Cạnh ngắn nhất là $MN = 15\text{ cm}$ (không phải QM).
    \item Ý d) Đúng: $15 + 25 = 40\text{ cm}$. Đúng.
\end{itemize}
\begin{itemize}
    \item a) \textbf{Đúng}.
    \item b) \textbf{Đúng}.
    \item c) \textbf{Sai}.
    \item d) \textbf{Đúng}.
\end{itemize}
}
\end{ex}

\textbf{Phần 3. Trả lời ngắn}
\setcounterpageref{ex}{0}

\begin{ex} %[2H5V26]
Đường gấp khúc gồm 3 đoạn thẳng có độ dài lần lượt là 18 cm, 25 cm và 37 cm. Tính độ dài đường gấp khúc đó.
\shortans{80}
\loigiai{
\textbf{Đáp số:} $80$\\
$18 + 25 + 37 = 80\text{ cm}$.
}
\end{ex}

\begin{ex} %[2H5V26]
Một hình tứ giác có chu vi là 75 cm, biết tổng độ dài ba cạnh là 56 cm. Tính độ dài cạnh còn lại của hình tứ giác.
\shortans{19}
\loigiai{
\textbf{Đáp số:} $19$\\
$75 - 56 = 19\text{ cm}$.
}
\end{ex}

\begin{ex} %[2H5V26]
Đường gấp khúc ABCD dài 48 cm gồm 3 đoạn thẳng có độ dài bằng nhau. Tính độ dài mỗi đoạn thẳng.
\shortans{16}
\loigiai{
\textbf{Đáp số:} $16$\\
$48 : 3 = 16\text{ cm}$.
}
\end{ex}

\begin{ex} %[2H5V26]
Tính chu vi hình tứ giác có độ dài các cạnh là 14 cm, 16 cm, 24 cm, 26 cm.
\shortans{80}
\loigiai{
\textbf{Đáp số:} $80$\\
$(14 + 16) + (24 + 26) = 30 + 50 = 80\text{ cm}$.
}
\end{ex}

\begin{ex} %[2H5V26]
Một đoạn dây dài 36 cm được uốn thành một hình tam giác có 3 cạnh bằng nhau. Hỏi mỗi cạnh dài bao nhiêu xăng-ti-mét?
\shortans{12}
\loigiai{
\textbf{Đáp số:} $12$\\
$36 : 3 = 12\text{ cm}$.
}
\end{ex}

\begin{ex} %[2H5V26]
Đường gấp khúc MNP gồm hai đoạn thẳng MN và NP có tổng độ dài là 42 cm. Đoạn thẳng MN dài 25 cm. Tính độ dài đoạn thẳng NP.
\shortans{17}
\loigiai{
\textbf{Đáp số:} $17$\\
$42 - 25 = 17\text{ cm}$.
}
\end{ex}
